\documentclass[10pt,a4paper]{amsart}
\usepackage[margin=19mm]{geometry}
\usepackage{amsmath,amssymb,mathtools}
\usepackage[scr=rsfs,cal=euler]{mathalfa}
\usepackage{xfrac}
\usepackage[unicode,hidelinks,bookmarks=false]{hyperref}
\newcommand{\ie}{i.e.\ }
\newcommand{\cf}{cf.\ }
\newcommand{\resp}{respectively\ }
\newcommand{\Subsection}{\subsection}
\providecommand{\nobreakdash}{\nobreak\hbox{-}}
\setlength{\emergencystretch}{2em}
\multlinegap0pt
\usepackage{bm}
\usepackage{bbm}
\usepackage{dsfont}
\usepackage{xspace}
\usepackage{cancel}
\usepackage{mathtools}
\usepackage{amsthm}
\makeatletter
\@ifundefined{theorem}{\newtheorem{theorem}{\bf Theorem}}{}
\makeatother
\title[Nitsche's method for the stationary Boussinesq system under ]{Nitsche's method for the stationary Boussinesq system under mixed and nonlinear boundary conditions}
\author{Aparna Bansal, Nicol\'as A. Barnafi, Gianmarco Sperone, Dwijendra N. Pandey}
\date{Source: arXiv:2604.06560v1 (CC BY 4.0)}
\begin{document}
\maketitle

\noindent\emph{Source and licence.} Nitsche's method for the stationary Boussinesq system under mixed and nonlinear boundary conditions. Version arXiv:2604.06560v1 (CC BY 4.0), distributed under the Creative Commons Attribution 4.0 International licence, as recorded in the arXiv licence metadata for this exact version and shown on the abstract page https://arxiv.org/abs/2604.06560v1. Full rights notice: licenses/arxiv-2604.06560-CC-BY-4.0.txt.\par\smallskip

\noindent\textit{Editorial scope.} Counted unit: the Theorem whose source label is \texttt{RR1} in the article (source0.tex lines 490--499, source label \texttt{RR1}), delivered together with the article's own complete original proof of it (source0.tex lines 500--566, one \texttt{proof} environment of 67 lines). No mathematics is written, repaired or re-derived: statement and proof are the article's own text line for line. Required context retained uncounted: bb (lines 117--119); E (lines 176--184); qa4 (lines 215--229); aa (lines 323--325); 51 (lines 326--332); 37 (lines 404--406). Every definition, display and label of those lines is retained unchanged. Omitted from this package and not redistributed: 1--116, 120--175, 185--214, 230--322, 333--403, 407--489, 567--1247. Nothing inside a retained excerpt is omitted or altered: every retained body is delivered exactly as the article's lines are written, so no omission receipt of any kind accompanies this package. Citations: the retained excerpts carry no citation command at all, so no bibliography entry of the article is transcribed and no reference is redistributed. Reference adaptation: none was needed and none was made; every reference inside the delivered unit resolves inside it, so no unresolved reference or citation is produced. Printed slips kept literal: no printed word, symbol, hypothesis or number of the counted statement or of its proof was altered. Layout and class: the article's own class is \texttt{article} with packages including latexsym, amssymb, bm, amsmath, mathtools, graphicx, epstopdf, epsfig, wrapfig, tikz, amsthm, float, caption, subcaption; the wrapper is the lane's pinned \texttt{amsart} editorial wrapper at 10pt on A4, which restates the theorem-like environments the retained text needs and adds the article's own macro definitions that the retained text needs (none), with extra wrapper packages (bm, bbm, dsfont, xspace, cancel, mathtools). The delivered numbering is the unit's own. Assets: no retained span contains a figure environment, a graphics-inclusion command, a PDF-page inclusion, a table or a TikZ picture; no image file is copied into this package, the package declares no asset dependency and no image file is redistributed with it. Licence: version arXiv:2604.06560v1 of this article is distributed under the Creative Commons Attribution 4.0 International licence, as recorded in the arXiv licence metadata for this exact version and shown on the abstract page https://arxiv.org/abs/2604.06560v1. A full rights notice is delivered at licenses/arxiv-2604.06560-CC-BY-4.0.txt. Characters: every retained excerpt is pure ASCII, so the delivered engine, XeLaTeX with its default Latin Modern text font, reports no missing character.
\begin{align}\label{bb}
	\mathcal{A}\left[\left(\boldsymbol{u}, p, \theta \right) ;\left(\boldsymbol {v}, q, \varphi \right)\right]=\mathcal{F}(\boldsymbol{v}, q, \varphi )  \qquad \forall  \left(\boldsymbol{v}, q, \varphi  \right) \in \mathrm{V}_{0} \times \Pi \times \mathrm{M}_{0},
	\end{align}

  	    \begin{equation}\label{E}
  	    \begin{aligned}
  	    \begin{array}{rlrl}
  	    {A}_S^h(\boldsymbol{u}_h, \boldsymbol{v}_h) + {O}_S^h(\boldsymbol{u}_h ; \boldsymbol{u}_h, \boldsymbol{v}_h) + {B}^h(\boldsymbol{v}_h,p_h) - D(\theta_h,\boldsymbol{v}_h)& = F_{v}(\boldsymbol{v}_h) & \forall \boldsymbol{v}_h \in \mathrm{V}_h, \\
  	    {B}^h(\boldsymbol{u}_h, q_h) & = F_q(q_h) & \forall q_h \in \mathrm{\Pi}_h, \\ 
       {A}^h_T(\theta_h,\varphi_h) + O_T(\boldsymbol{u}_h; \theta_h, \varphi_h) & = F_{\theta}(\varphi_h) & \forall \varphi_h \in \mathrm{M}_{h},
  	    \end{array}
  	    \end{aligned}
  	    \end{equation} 

  		\begin{theorem}\label{qa4}
  For a given function $\psi  \in L^{\infty}(\mathbb{R};\mathbb{R})$, there exist positive constants $C_1$, $C_2$, $C_3$, $C_4$, $C_5$, $C_6$, independent of $h$, such that
  		\begin{equation*}
  		\begin{aligned}
  		\left|{A}_S^h(\boldsymbol{u}_h, \boldsymbol{v}_h)\right| \leq& C_1 \|\boldsymbol{u}_h\|_{1, h}\|\boldsymbol{v}_h\|_{1,h}, &&  \forall \boldsymbol{u}_h, \boldsymbol{v}_h \in \mathrm{V}_h,
  		\\
  		\left|O_S^h\left(\boldsymbol{w}_h; \boldsymbol{u}_h, \boldsymbol{v}_h\right)\right| \leq& {C}_2 \left\|\boldsymbol{w}_h\right\|_{1,h}\|\boldsymbol{u}_h\|_{1,h}\|\boldsymbol{v}_h\|_{1, h}, && \forall \boldsymbol{w}_h, \boldsymbol{u}_h, \boldsymbol{v}_h \in \mathrm{V}_h,
  		\\
  		\left|{B}^h(\boldsymbol{v}_h, q_h)\right| \leq& C_3 \|\boldsymbol{v}_h\|_{1, h}\|q_h\|_{0, \Omega}, && \forall \boldsymbol{v}_h \in \mathrm{V}_h, q_h \in \Pi_h, \\ 
        \left| D( \theta_h, \boldsymbol{v}_h) \right|  \leq & C_4 \|\theta_h\|_{1,h}\|\boldsymbol{v}_h\|_{1,h} ,  && \forall \theta_h \in \mathrm{M}_h, \boldsymbol{v}_h \in \mathrm{V}_h, \\
        \left| A_T^h( \theta_h, \varphi_h)\right| \leq & C_5 \| \theta_h \|_{1,h} \| \varphi \|_{1,h} , && \forall  \theta_h, \varphi_h \in \mathrm{M}_h, \\
        \left| O_T^h(\boldsymbol{v}_h; \theta_h, \varphi_h) \right| \leq & C_6 \left\|\boldsymbol{v}_h\right\|_{1,h}\|\theta_h \|_{1,h}\|\varphi_h\|_{1,h} , && \forall \boldsymbol{v}_h \in \mathrm{V}_h, \theta_h, \varphi_h \in \mathrm{M}_h. 
  		\end{aligned}
  		\end{equation*}
        \end{theorem}  

  		\begin{align}\label{aa}
  		\mathcal{C}_h\left[(\boldsymbol{u}_h, \theta_h, p_h);( \boldsymbol{v}_h,\varphi_h, q_h)\right]={A}_S^h(\boldsymbol{u}_h, \boldsymbol{v}_h) + {B}^h(\boldsymbol{v}_h,p_h) + {B}^h(\boldsymbol{u}_h, q_h) +  {A}^h_T(\theta_h,\varphi_h) .
  		\end{align}

  		\begin{theorem}\label{51}
  			There exist a positive constant $\hat{\alpha}= \hat{\alpha}(C_S,C_T,\hat{\theta})$ such that 
  			\begin{align*}
  			\sup _{\boldsymbol{0} \neq (\boldsymbol{v}_h,\varphi_h , q_h) \in \mathrm{V}_h \times \mathrm{M}_h \times  \mathrm{\Pi}_h}\frac{\mathcal{C}_h\left[\left( \boldsymbol{u}_h, \theta_h, p_h\right);\left( \boldsymbol{v}_h, \varphi_h, q_h \right)\right]}{\left\|\left( \boldsymbol{v}_h, \varphi_h, q_h \right)\right\|} \geq \hat{\alpha}\left\|\left(\boldsymbol{u}_h,\theta_h, p_h\right)\right\| \quad \forall\left( \boldsymbol{u}_h,\varphi_h,p_h\right) \in \mathrm{V}_h \times \mathrm{M}_h \times \mathrm{\Pi}_h,
  			\end{align*}
  			with $C_S$, $C_T$ and $\hat{\theta}$ are the coercivity and inf-sup stability constants. 
  		\end{theorem}

  			\begin{align}\label{37}
  			\sup _{\boldsymbol{0} \neq (\boldsymbol{v}_h, \varphi_h, q_h) \in \mathrm{V}_h \times \mathrm{M}_h \times \Pi_h} \mathcal{A}_{\boldsymbol{w}_h}\left[(\boldsymbol{v}_h, \varphi_h, q_h);(\boldsymbol{u}_h, \theta_h, p_h)\right] & \geq \frac{\hat{\alpha}}{2}\|(\boldsymbol{u}_h, \theta_h, p_h)\| >0  \\ & \forall(\boldsymbol{u}_h,\theta_h, p_h) \in \mathrm{V}_h \times \mathrm{M}_h \times \Pi_h,(\boldsymbol{u}_h, \theta_h, p_h) \neq 0.
  			\end{align}

  		\begin{theorem}\label{RR1}
  			Assume that
  			\begin{align}\label{57}
  			\| \boldsymbol{u} \|_{1}    + \| \theta \|_{1} \leq \frac{\hat{\alpha}}{4},
  			\end{align}
  			with $\hat{\alpha}$ being the positive constant in Theorem~\ref{51}. Let $(\boldsymbol{u},\theta, p) \in$ $\mathrm{V} \times \mathrm{M} \times \Pi$ and $\left(\boldsymbol{u}_h,\theta_h,p_h \right) \in \mathrm{V}_h \times \mathrm{M}_h \times\mathrm{\Pi}_h$ be the unique solutions of problems \eqref{bb} and \eqref{E}, respectively. Then, we have 
  			\begin{align}\label{58}
  			\left\|\left( \boldsymbol{u}-\boldsymbol{u}_h, \theta - \theta_h, p - p_h \right)\right\| \lesssim  \inf _{\boldsymbol{0} \neq  \left(\boldsymbol{v}_h,\varphi_h,q_h \right) \in \mathrm{V}_h \times \mathrm{M}_h \times  \mathrm{\Pi}_h}\left\|\left(\boldsymbol{u}-\boldsymbol{v}_h, \varphi - \varphi_h, p-q_h \right)\right\|.
  			\end{align}	
  		\end{theorem} 

  		\begin{proof}
  			In order to simplify the subsequent analysis, we define $\boldsymbol{e}_{\mathrm{u}} = \boldsymbol{u}-\boldsymbol{u}_h, \boldsymbol{e}_{\theta} = \theta -\theta_h,$ and  $\boldsymbol{e}_p = p - p_h$, and for any $\left( \boldsymbol{z}_h, \phi_h, \zeta_h \right) \in \mathrm{V}_h \times \mathrm{M}_h \times \mathrm{\Pi}_h $, we write
  			\begin{align}\label{error_d}
  			\boldsymbol{e}_{\boldsymbol{u}}&=\xi_{\boldsymbol{u}}+\chi_{\boldsymbol{u}}=\left(\boldsymbol{u}-\boldsymbol{z}_h\right)+\left(\boldsymbol{z}_h-\boldsymbol{u}_h\right), \nonumber \\
            \boldsymbol{e}_{\theta}&=\xi_{\theta}+\chi_{\theta}=\left(\theta -\phi_h\right)+\left(\phi_h-\theta _h\right), 
            \\   \boldsymbol{e}_p&=\xi_p+\chi_p=\left(p-\zeta_h\right)+\left(\zeta_h-p_h\right) \nonumber .
  			\end{align}
  			By recalling the definition of the bilinear forms, we observe the validity of following identities:
  			$$
  			\mathcal{C}_h\left[( \boldsymbol{u},\theta,p);(\boldsymbol{v}_h,\varphi_h,q_h)\right] - D(\theta, \boldsymbol{v}_h)+O_h^S(\boldsymbol{u} ; \boldsymbol{u}, \boldsymbol{v}_h) +O_T( \boldsymbol{u}; \theta, \varphi_h)= \mathcal{F}(\boldsymbol{v}_h, \varphi_h, q_h),
  			$$
  			and
  			$$
  			\mathcal{C}_h\left[( \boldsymbol{u}_h,\theta_h,p_h);(\boldsymbol{v}_h,\varphi_h,q_h)\right] - D(\theta_h, \boldsymbol{v}_h)+O_h^S(\boldsymbol{u}_h ; \boldsymbol{u}_h, \boldsymbol{v}_h) +O_T( \boldsymbol{u}_h; \theta_h, \varphi_h)= \mathcal{F}(\boldsymbol{v}_h, \varphi_h, q_h), 
  			$$
  			for all $(\boldsymbol{v}_h,\varphi_h,q_h) \in \mathrm{V}_h \times \mathrm{M}_h \times \Pi_h$. Based on these observations, we deduce the Galerkin orthogonality property
     \begin{align}\label{galerkin_ortho}
\mathcal{C}_h\left[(\boldsymbol{e}_{\boldsymbol{u}}, \boldsymbol{e}_{\theta}, \boldsymbol{e}_p);( \boldsymbol{v}_h, \varphi_h, q_h)\right]  - D (\boldsymbol{e}_{\theta}, \boldsymbol{v}_h) + \left[O_S^h\left(\boldsymbol{u}; \boldsymbol{u}, \boldsymbol{v}_h \right) -O_
S^h \left(\boldsymbol{u}_h; \boldsymbol{u}_h, \boldsymbol{v}_h \right)\right] + \left[O_T\left(\boldsymbol{u}; \theta , \varphi_h \right) \nonumber \right. \\ \quad \left.  - O_T \left(\boldsymbol{u}_h; \theta_h, \varphi_h \right)\right] = 0,  \forall \left( \boldsymbol{v}_h, \varphi_h, q_h \right) \in \mathrm{V}_h \times \mathrm{M}_h \times {\Pi}_h.
\end{align}
     Now, using  \eqref{error_d}, and \eqref{galerkin_ortho}, we deduce that for all $\left( \boldsymbol{v}_h,\varphi_h, q_h \right) \in \mathrm{V}_h \times \mathrm{M}_h \times \mathrm{\Pi}_h $, the following relationship holds
     $$
     \begin{aligned}
& \mathcal{A}_{\boldsymbol{u}_{\mathrm{h}}}\left[\left( \chi_{\boldsymbol{u}},\chi_{\theta},\chi_p\right);\left(\boldsymbol{v}_h,\varphi_h,q_h \right)\right] - D(\chi_{\theta}, \boldsymbol{v}_h) \\ 
& \hspace{-4mm} = 
  \mathcal{C}_h \left[\left( \chi_{\boldsymbol{u}}, \chi_{\theta}, \chi_p \right), \left( \boldsymbol{v}_h, \varphi_h, q_h \right)\right] - D(\chi_{\theta}, \boldsymbol{v}_h) + O_S^h\left(\boldsymbol{u}_h; \chi_{\boldsymbol{u}}, \boldsymbol{v}_h\right) + O_T(\boldsymbol{u}_h; \chi_{\theta}, \varphi_h ) \\
  & \hspace{-4mm} = - \mathcal{C}_h\left[\left( \xi_{\boldsymbol{u}}, \xi_{\theta}, \xi_p \right); \left( \boldsymbol{v}_h, \varphi_h, q_h\right)\right] +
  \mathcal{C}_h\left[\left( \boldsymbol{e}_{\boldsymbol{u}}, \boldsymbol{e}_{\theta}, \boldsymbol{e}_p \right); \left( \boldsymbol{v}_h, \varphi_h, q_h\right)\right] + D(\xi_{\theta}, \boldsymbol{v}_h) - D(\boldsymbol{e}_{\theta}, \boldsymbol{v}_h) \\
  & + O_S^h\left(\boldsymbol{u}_h; \chi_{\boldsymbol{u}}, \boldsymbol{v}_h\right) + O_T(\boldsymbol{u}_h; \chi_{\theta}, \varphi_h ) \\
  & \hspace{-4mm} =- \mathcal{C}_h\left[\left( \xi_{\boldsymbol{u}}, \xi_{\theta}, \xi_p \right); \left( \boldsymbol{v}_h, \varphi_h,q_h\right)\right] + D(\xi_{\theta}, \boldsymbol{v}_h) - \left[O_S^h\left(\boldsymbol{u}; \boldsymbol{u}, \boldsymbol{v}_h \right) -O_
S^h \left(\boldsymbol{u}_h; \boldsymbol{u}_h, \boldsymbol{v}_h \right)\right] \\ 
& - \left[O_T\left(\boldsymbol{u}; \theta , \varphi_h \right) - O_T \left(\boldsymbol{u}_h; \theta_h, \varphi_h \right)\right] + O_S^h\left(\boldsymbol{u}_h; \chi_{\boldsymbol{u}}, \boldsymbol{v}_h\right) + O_T(\boldsymbol{u}_h; \chi_{\theta}, \varphi_h ) \\
  & \hspace{-4mm} =- \mathcal{C}_h\left[\left( \xi_{\boldsymbol{u}}, \xi_{\theta}, \xi_p \right); \left( \boldsymbol{v}_h, \varphi_h,q_h\right)\right] + D(\xi_{\theta}, \boldsymbol{v}_h) - \left[O_S^h\left(\boldsymbol{u} - \boldsymbol{u}_h; \boldsymbol{u}, \boldsymbol{v}_h \right) \right.  \\ & \quad \left. + O_S^h\left( \boldsymbol{u}_h; \boldsymbol{u}, \boldsymbol{v}_h\right) -O_
S^h \left(\boldsymbol{u}_h; \boldsymbol{u}_h, \boldsymbol{v}_h \right)\right]  - \left[O_T\left(\boldsymbol{u} - \boldsymbol{u}_h; \theta , \varphi_h \right) \right.  \\ & \quad \left. +O_T\left(\boldsymbol{u}_h; \theta, \varphi_h \right)  - O_T \left(\boldsymbol{u}_h; \theta_h, \varphi_h \right)\right] + O_S^h\left(\boldsymbol{u}_h; \chi_{\boldsymbol{u}}, \boldsymbol{v}_h\right) + O_T(\boldsymbol{u}_h; \chi_{\theta}, \varphi_h )  \\
  & \hspace{-4mm} =- \mathcal{C}_h\left[\left( \xi_{\boldsymbol{u}}, \xi_{\theta}, \xi_p \right); \left( \boldsymbol{v}_h, \varphi_h,q_h\right)\right] + D(\xi_{\theta}, \boldsymbol{v}_h) - O_S^h\left(\xi_{\boldsymbol{u}} + \chi_{\boldsymbol{u}}; \boldsymbol{u}, \boldsymbol{v}_h \right)  \\ & \quad - O_S^h\left( \boldsymbol{u}_h; \xi_{\boldsymbol{u}} + \chi_{\boldsymbol{u}}, \boldsymbol{v}_h\right) + O_S^h\left(\boldsymbol{u}_h; \chi_{\boldsymbol{u}}, \boldsymbol{v}_h\right) -  O_T\left(\xi_{\boldsymbol{u}} + \chi_{\boldsymbol{u}}; \theta , \varphi_h \right) \\& \quad - O_T\left(\boldsymbol{u}_h; \xi_{\theta} + \chi_{\theta}, \varphi_h \right) + O_T(\boldsymbol{u}_h; \chi_{\theta}, \varphi_h )  \\
  & \hspace{-4mm} =- \mathcal{C}_h\left[\left( \xi_{\boldsymbol{u}}, \xi_{\theta}, \xi_p \right); \left( \boldsymbol{v}_h, \varphi_h,q_h\right)\right] + D(\xi_{\theta}, \boldsymbol{v}_h) - O_S^h\left(\xi_{\boldsymbol{u}}; \boldsymbol{u}, \boldsymbol{v}_h \right) - O_S^h\left(\chi_{\boldsymbol{u}}; \boldsymbol{u}, \boldsymbol{v}_h \right) \\ & \quad   - O_S^h\left( \boldsymbol{u}_h; \xi_{\boldsymbol{u}}, \boldsymbol{v}_h\right) -  O_T\left( \xi_{\boldsymbol{u}}; \theta , \varphi_h \right) -  O_T\left( \chi_{\boldsymbol{u}}; \theta , \varphi_h \right)  - O_T\left(\boldsymbol{u}_h;  \xi_{\theta}, \varphi_h \right)
\end{aligned}
$$
 which, together with the definition of $\mathcal{C}_h$ given in equation \eqref{aa}, implies
\begin{align}\label{61}
  \mathcal{A}_{\boldsymbol{u}_{\mathrm{h}}}\left[\left( \chi_{\boldsymbol{u}}, \chi_{\theta}, \chi_p \right); \left( \boldsymbol{v}_h, \varphi_h, q_h \right)\right] - D(\chi_{\theta}, \boldsymbol{v}_h)  & = -{A}_S^h(\xi_{\boldsymbol{u}}, \boldsymbol{v}_h) - {B}^h(\boldsymbol{v}_h,\xi_{p}) - {B}^h(\xi_{\boldsymbol{u}}, q_h) -  {A}^h_T(\xi_{\theta},\varphi_h) + D(\xi_{\theta}, \boldsymbol{v}_h) \nonumber \\ & \quad  - O_S^h\left(\xi_{\boldsymbol{u}}; \boldsymbol{u}, \boldsymbol{v}_h \right) - O_S^h\left(\chi_{\boldsymbol{u}}; \boldsymbol{u}, \boldsymbol{v}_h \right) - O_S^h\left( \boldsymbol{u}_h; \xi_{\boldsymbol{u}}, \boldsymbol{v}_h\right) \nonumber  \\ & \quad -  O_T\left( \xi_{\boldsymbol{u}}; \theta , \varphi_h \right)   -  O_T\left( \chi_{\boldsymbol{u}}; \theta , \varphi_h \right) - O_T\left(\boldsymbol{u}_h;  \xi_{\theta}, \varphi_h \right), 
\end{align}
for all $\left( \boldsymbol{v}_h,\varphi_h,q_h\right) \in \mathrm{V}_h \times \mathrm{M}_h \times \mathrm{\Pi}_h$. Next, utilizing the discrete inf-sup condition \eqref{37} at the left hand side of \eqref{61}, and applying the continuity properties of $A_S^h, A_T^h, B^h, D, O_T, O_S^h$ stated in Theorem~\ref{qa4} to the right hand side of \eqref{61}, we can derive
   \begin{align*}
  & \left\|\chi_{\boldsymbol{u}}\right\|_{1,h} +   \left\|\chi_{\theta}\right\|_{1,h} + \left\|\chi_p\right\|_{0} \\
  & \hspace{-4mm} \lesssim \frac{2}{\hat{\alpha} \| (\boldsymbol{v}_h, \varphi_h, q_h) \|} \bigg( \|\xi_{\boldsymbol{u}}\|_{1,h} \|\boldsymbol{v}_h\|_{1,h} + \|\boldsymbol{v}_h\|_{1,h} \| \xi_p\|_{0} + \|\xi_{\boldsymbol{u}}\|_{1,h} \|q_h\|_{0} + \| \xi_{\theta}\|_{1,h} \| \varphi_h \|_{1,h} \\& \quad + \| \xi_{\theta} \|_{1,h}   \|\boldsymbol{v}_h\|_{1,h}  + \|\xi_{\boldsymbol{u}}\|_{1,h} \|\boldsymbol{v}_h\|_{1,h} \| \boldsymbol{u} \|_{1}  + \|\chi_{\boldsymbol{u}}\|_{1,h} \|\boldsymbol{v}_h\|_{1,h} \| \boldsymbol{u} \|_{1}  + \|\xi_{\boldsymbol{u}}\|_{1,h} \|\boldsymbol{v}_h\|_{1,h} \| \boldsymbol{u}_h \|_{1,h} \\
  &\quad  + \| \xi_{\boldsymbol{u}} \|_{1,h} \| \theta \|_{1} \| \varphi_h \|_{1,h} + \| \chi_{\boldsymbol{u}} \|_{1,h} \| \theta \|_{1} \| \varphi_h \|_{1,h} + \| \boldsymbol{u}_h \|_{1,h} \| \xi_{\theta} \|_{1,h} \| \varphi_h \|_{1,h} \bigg)  \\      
  & \hspace{-4mm} \lesssim \frac{2}{\hat{\alpha}} \bigg( \|\xi_{\boldsymbol{u}}\|_{1,h}  +  \| \xi_p\|_{0,\Omega} + \|\xi_{\boldsymbol{u}}\|_{1,h} + \| \xi_{\theta}\|_{1,h} + \| \xi_{\theta} \|_{1,h}  + \|\xi_{\boldsymbol{u}}\|_{1,h} \| \boldsymbol{u} \|_{1} + \|\chi_{\boldsymbol{u}}\|_{1,h} \| \boldsymbol{u} \|_{1}  \\
  &\quad  + \|\xi_{\boldsymbol{u}}\|_{1,h} \| \boldsymbol{u}_h \|_{1,h}  + \| \xi_{\boldsymbol{u}} \|_{1,h} \| \theta \|_{1,h}  + \| \chi_{\boldsymbol{u}} \|_{1,h} \| \theta \|_{1} + \| \boldsymbol{u}_h \|_{1} \| \xi_{\theta} \|_{1,h}  \bigg) \\
  & \hspace{-4mm} \lesssim \frac{2}{\hat{\alpha}} \bigg(  \| \xi_p\|_{0,\Omega} + \| \xi_{\theta} \|_{1,h} + \|\xi_{\boldsymbol{u}}\|_{1,h} \left( 1+ \| \boldsymbol{u} \|_{1} + \| \boldsymbol{u}_h \|_{1,h} +  \| \theta \|_{1,h} \right)  \\
  &\quad + \|\chi_{\boldsymbol{u}}\|_{1,h} \left(\| \boldsymbol{u} \|_{1}    + \| \theta \|_{1} \right) + \| \boldsymbol{u}_h \|_{1,h} \| \xi_{\theta} \|_{1,h}  \bigg),
\end{align*}
as well as
 \begin{align*}
& \left( 1- \frac{2}{\hat{\alpha}} \left(\| \boldsymbol{u} \|_{1}    + \| \theta \|_{1} \right) \right)  \left\|\chi_{\boldsymbol{u}}\right\|_{1,h} +  \left\|\chi_{\theta}\right\|_{1,h} + \left\|\chi_p\right\|_{0} \\
& \hspace{-4mm} \lesssim \frac{2}{\hat{\alpha}} \bigg(  \| \xi_p\|_{0,\Omega} + \| \xi_{\theta} \|_{1,h} +  \left( 1+ \| \boldsymbol{u} \|_{1} + \| \boldsymbol{u}_h \|_{1,h} +  \| \theta \|_{1,h} \right) \|\xi_{\boldsymbol{u}}\|_{1,h} + \| \boldsymbol{u}_h \|_{1,h} \| \xi_{\theta} \|_{1,h}  \bigg).
\end{align*}
  			Therefore, by taking into account the assumption \eqref{57} and the fact that $(\boldsymbol{u}_h, \theta_h) \in \boldsymbol{K}_h$, we can conclude
  			\begin{align}\label{63}
  			\left\|\chi_p\right\|_{0}+\left\|\chi_{\boldsymbol{u}}\right\|_{1,h} +\left\|\chi_{\theta}\right\|_{1,h} \lesssim \left(\left\|\xi_p\right\|_{0}+\left\|\xi_{\boldsymbol{u}}\right\|_{1,h} + \left\| \xi_{\theta}\right\|_{1,h}\right).
  			\end{align}
            In this way, from \eqref{error_d}, \eqref{63} and the triangle inequality we obtain
  			$$
  			\left\|\left(\boldsymbol{e}_p, \boldsymbol{e}_{\boldsymbol{u}},  \boldsymbol{e}_{\theta} \right)\right\| \leq\left\|\left(\chi_p, \chi_{\boldsymbol{u}}, \chi_{\theta}\right)\right\|+\left\|\left({\xi}_p, {\xi}_{\boldsymbol{u}}, \xi_{\theta}\right)\right\| \lesssim \left\|\left({\xi}_p, {\xi}_{\mathrm{u}}, \xi_{\theta}\right)\right\|.
  			$$
  		This, together with the fact that  $\left(\boldsymbol{z}_h,\zeta_h, \phi_h \right) \in \mathrm{V}_h \times \mathrm{\Pi}_h \times \mathrm{M}_h $ is arbitrary, concludes the proof.
  		\end{proof}

\end{document}
